\chapter{Test Codes and Ordering --- Talking to Your Doctor}
\label{app:test_codes}

\clearpage
\section{Introduction}

This appendix is designed for patients who want to have informed conversations with their doctors about which tests to order. It includes examples of LOINC codes, CPT codes, and blood tube colors. Code systems are versioned and updated; the examples should be checked against the current LOINC release, CPT code set, payer rules, and laboratory ordering catalog.

You do not need to memorize these codes. But if your doctor says ``we will order a CBC,'' knowing that ``CBC'' is CPT code 85025 and means a complete blood count with differential gives you a better understanding of what is actually being tested and what you are paying for.

\clearpage
\section{Common Blood Tests --- LOINC Codes}

LOINC (Logical Observation Identifiers Names and Codes) is an international terminology for identifying observations, including laboratory tests. A LOINC code does not by itself guarantee identical specimen, method, reference interval, reporting units, or ordering availability across laboratories. CPT is a US procedure and billing code set and should not be treated as a universal test identifier.

For maintenance, record the LOINC release date and the CPT code-set year used for this appendix. Confirm codes with the current official LOINC database, AMA CPT code set, payer policy, and the performing laboratory before using them for ordering or billing.

\begin{longtable}{@{}p{5.5cm}p{2.5cm}p{2.2cm}@{}}
\toprule
\textbf{Test Name} & \textbf{LOINC Code} & \textbf{Specimen} \\
\midrule
\endhead
\textbf{Complete Blood Count (CBC)} & & \\
CBC with differential & 58410-2 & Blood \\
Hemoglobin & 718-7 & Blood \\
Hematocrit & 4544-3 & Blood \\
WBC count & 6690-2 & Blood \\
Platelet count & 777-3 & Blood \\
RBC count & 789-8 & Blood \\
MCV & 785-6 & Blood \\
MCH & 786-4 & Blood \\
MCHC & 787-2 & Blood \\
RDW & 788-0 & Blood \\
Reticulocyte count & 10405-5 & Blood \\
\textbf{Inflammation} & & \\
ESR (sed rate) & 4537-7 & Blood \\
CRP (C-reactive protein) & 30522-7 & Serum \\
\textbf{Coagulation} & & \\
PT (prothrombin time) & 5902-2 & Plasma \\
INR & 38885-1 & Plasma \\
aPTT & 14959-1 & Plasma \\
Fibrinogen & 3255-7 & Plasma \\
D-dimer & 7738-4 & Plasma \\
\textbf{Metabolic Panel} & & \\
Glucose & 2345-7 & Serum/Plasma \\
BUN & 3094-0 & Serum \\
Creatinine & 2160-0 & Serum \\
eGFR & 48642-3 & Serum \\
Sodium & 2951-2 & Serum \\
Potassium & 2823-3 & Serum \\
Chloride & 2075-0 & Serum \\
CO$_2$ (bicarbonate) & 2028-9 & Serum \\
Calcium & 17861-6 & Serum \\
Magnesium & 19221-5 & Serum \\
Phosphorus & 2777-1 & Serum \\
\textbf{Liver} & & \\
Total protein & 2885-2 & Serum \\
Albumin & 1751-7 & Serum \\
Total bilirubin & 1975-2 & Serum \\
Direct bilirubin & 1968-7 & Serum \\
AST & 1920-8 & Serum \\
ALT & 1742-6 & Serum \\
Alkaline phosphatase & 6768-6 & Serum \\
GGT & 2324-4 & Serum \\
LDH & 14828-2 & Serum \\
\textbf{Cardiac} & & \\
Troponin I & 10839-9 & Serum \\
Troponin T & 49563-0 & Serum \\
BNP & 30934-4 & Plasma \\
NT-proBNP & 33762-4 & Plasma \\
\textbf{Endocrine} & & \\
HbA1c & 4548-4 & Blood \\
TSH & 3016-3 & Serum \\
Free T4 & 3026-2 & Serum \\
Free T3 & 3027-0 & Serum \\
Cortisol & 2143-6 & Serum \\
\textbf{Uric Acid} & & \\
Uric acid & 14957-5 & Serum \\
\textbf{Blood Gas} & & \\
Calcium, ionized & 17144-7 & Blood \\
Lactate & 2518-9 & Blood \\
pH (blood gas) & 2744-1 & Blood \\
pCO$_2$ & 2021-4 & Blood \\
pO$_2$ & 2703-7 & Blood \\
\textbf{Tumor Markers} & & \\
PSA & 2857-1 & Serum \\
\textbf{Urinalysis} & & \\
Urinalysis (complete) & 24331-1 & Urine \\
Urine protein & 2888-6 & Urine \\
Urine creatinine & 2161-8 & Urine \\
Urine glucose & 2346-5 & Urine \\
\bottomrule
\end{longtable}

\clearpage
\section{Common Blood Tests --- CPT Codes}

CPT (Current Procedural Terminology) codes are what your insurance company uses to bill for your lab tests. Knowing these codes helps you understand what you are being charged for and helps you compare costs if you are paying out of pocket or shopping for the best price.

\begin{longtable}{@{}p{5.5cm}p{2.5cm}p{2.5cm}@{}}
\toprule
\textbf{Test Name} & \textbf{CPT Code} & \textbf{Category} \\
\midrule
\endhead
\textbf{Hematology} & & \\
CBC with differential & 85025 & Hematology \\
CBC without differential & 85027 & Hematology \\
WBC count & 85048 & Hematology \\
Platelet count & 85049 & Hematology \\
Reticulocyte count & 85044 & Hematology \\
\textbf{Coagulation} & & \\
PT & 85610 & Coagulation \\
INR & 85610 & Coagulation \\
aPTT & 85730 & Coagulation \\
Fibrinogen & 85384 & Coagulation \\
D-dimer (quantitative) & 85379 & Coagulation \\
\textbf{Chemistry} & & \\
Glucose & 82947 & Chemistry \\
BUN & 84520 & Chemistry \\
Creatinine & 82565 & Chemistry \\
eGFR (calculated) & 82565 & Chemistry \\
Comprehensive Metabolic Panel (CMP) & 80053 & Chemistry \\
Basic Metabolic Panel (BMP) & 80048 & Chemistry \\
Sodium & 84295 & Chemistry \\
Potassium & 84132 & Chemistry \\
Chloride & 82435 & Chemistry \\
CO$_2$ & 82374 & Chemistry \\
Calcium & 82310 & Chemistry \\
Magnesium & 83735 & Chemistry \\
Phosphorus & 84100 & Chemistry \\
\textbf{Liver} & & \\
Total protein & 84155 & Chemistry \\
Albumin & 82040 & Chemistry \\
Total bilirubin & 82247 & Chemistry \\
Direct bilirubin & 82248 & Chemistry \\
AST & 84450 & Chemistry \\
ALT & 84460 & Chemistry \\
Alkaline phosphatase & 84075 & Chemistry \\
GGT & 82977 & Chemistry \\
LDH & 83615 & Chemistry \\
Hepatic Function Panel & 80076 & Chemistry \\
\textbf{Cardiac} & & \\
Troponin I & 84484 & Chemistry \\
BNP & 83880 & Chemistry \\
NT-proBNP & 83880 & Chemistry \\
\textbf{Endocrine} & & \\
HbA1c & 83036 & Chemistry \\
TSH & 84443 & Endocrinology \\
Free T4 & 84439 & Endocrinology \\
Free T3 & 84481 & Endocrinology \\
Cortisol & 82533 & Chemistry \\
\textbf{Tumor Markers} & & \\
PSA (total) & 86300 & Tumor marker \\
PSA (free) & 86316 & Tumor marker \\
\textbf{Immunology} & & \\
CRP (high-sensitivity) & 86141 & Immunology \\
ANA & 86038 & Immunology \\
Anti-dsDNA & 86200 & Immunology \\
RF (rheumatoid factor) & 86431 & Immunology \\
Anti-CCP & 86200 & Immunology \\
\textbf{Infectious Disease} & & \\
Hepatitis B surface Ag & 87340 & Microbiology \\
Hepatitis C Ab & 86704 & Microbiology \\
HIV Ag/Ab combo & 87389 & Microbiology \\
RPR & 86592 & Microbiology \\
Blood culture & 87040 & Microbiology \\
Urine culture & 87088 & Microbiology \\
Gram stain & 87205 & Microbiology \\
\textbf{Urinalysis} & & \\
Urinalysis (dipstick) & 81001 & Urinalysis \\
Urinalysis (microscopic) & 81015 & Urinalysis \\
Pregnancy test (urine) & 81025 & Urinalysis \\
HCG (serum, quantitative) & 84702 & Endocrinology \\
\bottomrule
\end{longtable}

\clearpage
\section{Blood Tube Colors --- What Is in Each Tube}

Every color-coded tube has a specific additive inside that prepares the blood for a particular type of test. The wrong tube gives the wrong results. This is what the phlebotomist is thinking about when they draw your blood.

\begin{longtable}{@{}p{3.2cm}p{4.2cm}p{5.5cm}@{}}
\toprule
\textbf{Tube Color} & \textbf{Additive Inside} & \textbf{Used For} \\
\midrule
Red (no additive) & None & Chemistry, serology (blood clot, serum separates out) \\
Gold/SST (red-gray top) & Gel separator + clot activator & Chemistry, endocrinology, antibodies \\
Green (dark green top) & Heparin (lithium or sodium) & Chemistry, blood gas \\
Light green (mint top) & Lithium heparin + gel separator & General chemistry \\
Lavender/Purple (purple top) & EDTA & CBC, HbA1c, blood typing, DNA tests \\
Blue (light blue top) & Sodium citrate & Coagulation tests (PT/INR, aPTT, D-dimer) \\
Gray (gray top) & Sodium fluoride + potassium oxalate & Glucose, lactate (prevents glycolysis) \\
Black (black top) & Sodium citrate (different ratio) & ESR (erythrocyte sedimentation rate) \\
Pink (pink top) & EDTA & Blood bank (type and crossmatch) \\
Yellow (yellow top) & ACD (acid citrate dextrose) & Molecular testing, DNA studies \\
\bottomrule
\end{longtable}

\subsubsection{What This Means For You}

\textbf{Why tube color matters:} If the wrong tube is used, the test may be inaccurate or uninterpretable. For example:
\begin{itemize}
    \item If a coagulation test (PT/INR) is drawn in the wrong tube, the result is unreliable and your warfarin dose could be adjusted incorrectly.
    \item If a glucose is drawn in a red-top tube instead of a gray-top tube, the glucose will be falsely low because the blood cells will consume the sugar during clotting.
    \item If a CBC is drawn in a green-top tube instead of a lavender-top tube, the anticoagulant will not preserve the cells properly.
\end{itemize}

You do not need to tell the phlebotomist which tube to use --- this is their job. But if you see them reach for the wrong tube, it is perfectly appropriate to ask: ``Is that the right tube for my test?''

\clearpage
\section{Panel Codes --- What You Are Actually Ordering}

Most lab tests are ordered as panels (groups of related tests) rather than individually. Here is what you are getting with each common panel.

\begin{longtable}{@{}p{4.0cm}p{3.0cm}p{6.0cm}@{}}
\toprule
\textbf{Panel Name} & \textbf{CPT Code} & \textbf{Tests Included} \\
\midrule
Basic Metabolic Panel (BMP) & 80048 & Glucose, BUN, creatinine, sodium, potassium, chloride, CO$_2$ \\
Comprehensive Metabolic Panel (CMP) & 80053 & BMP + calcium, total protein, albumin, bilirubin, AST, ALT, ALP \\
Hepatic Function Panel & 80076 & Bilirubin, albumin, AST, ALT, ALP, GGT \\
Lipid Panel & 80061 & Total cholesterol, LDL, HDL, triglycerides \\
Thyroid Panel & 84443 + 84439 & TSH + Free T4 (sometimes Free T3) \\
CBC with Differential & 85025 & WBC count, RBC, hemoglobin, hematocrit, platelets, indices, differential \\
Coagulation Panel & 85610 + 85730 & PT/INR + aPTT (often ordered together) \\
Iron Studies & 83540 + 83550 + 82728 & Iron, TIBC, ferritin (\% saturation is calculated) \\
Urinalysis (complete) & 81001 + 81015 & Dipstick + microscopic examination \\
\bottomrule
\end{longtable}

\clearpage
\section{How to Talk to Your Doctor About Lab Tests}

You do not need a medical degree to have a productive conversation about lab tests. Here is a framework.

\subsection{What Preparation Is Needed}

\begin{enumerate}
    \item \textbf{Write down what symptoms you have and when they started.} This helps your doctor decide which tests to order.
    \item \textbf{Make a list of all your medications, supplements, and vitamins.} This affects which tests are accurate and which may be interfered with.
    \item \textbf{Know your family medical history.} If your parents had diabetes, heart disease, or thyroid problems, you may need earlier screening.
    \item \textbf{Ask yourself what you want to know.} ``Do I want to check for diabetes?'' ``Do I want to understand why I am tired?'' ``Do I want to screen for heart disease risk?''
\end{enumerate}

\subsection{When a Test Is Ordered}

\begin{itemize}
    \item ``What test are you ordering and what will it tell us?''
    \item ``Do I need to fast for this test?''
    \item ``Are any of my medications likely to affect the results?''
    \item ``Should I stop any supplements before the test?''
    \item ``When will the results be available?''
    \item ``What happens if the results are abnormal?''
    \item ``How much will this test cost with my insurance?''
    \item ``Is there a less expensive alternative test that gives the same information?''
\end{itemize}

\subsection{Understanding Your Results}

\begin{itemize}
    \item ``Is this result normal for someone my age and sex?''
    \item ``How does this compare to my previous results?''
    \item ``Is this result trending in a concerning direction?''
    \item ``Do I need any follow-up testing?''
    \item ``Do I need treatment, or should we recheck this in a few months?''
    \item ``Could anything have affected this result (food, exercise, medication)?''
\end{itemize}

\clearpage
\section{Understanding Your Lab Bill}

Lab billing can be confusing. Here is what to know:

\begin{itemize}
    \item \textbf{Each test has a CPT code.} If you see multiple CPT codes on your bill, you are being charged for each test separately.
    \item \textbf{Panels are bundled.} A ``Comprehensive Metabolic Panel'' (CPT 80053) is one charge that includes 14 individual tests. You should not be billed separately for glucose, sodium, potassium, etc., if a CMP was ordered.
    \item \textbf{If you are paying out of pocket, ask for the cash price.} Lab cash prices are often dramatically lower than billed rates. Quest Diagnostics and LabCorp both publish self-pay pricing.
    \item \textbf{If your claim is denied, ask for an itemized bill} and verify that the correct CPT codes were used. Insurance denials are often due to coding errors.
    \item \textbf{Preventive screening tests are often covered at no cost} under the Affordable Care Act. This includes cholesterol screening, diabetes screening, and certain other tests as recommended by the USPSTF.
\end{itemize}
